\documentclass[11pt]{article}
\usepackage[margin=1in]{geometry}
\usepackage{amsmath,amssymb,amsthm}
\usepackage{booktabs}
\usepackage{hyperref}

\newtheorem{assumption}{Assumption}
\newtheorem{lemma}{Lemma}
\newtheorem{proposition}{Proposition}
\newtheorem{theorem}{Theorem}
\newtheorem{corollary}{Corollary}
\newcommand{\R}{\mathbb{R}}
\newcommand{\dist}{\operatorname{dist}}
\newcommand{\prox}{\operatorname{prox}}
\newcommand{\Fix}{\operatorname{Fix}}

\title{Local Sharpness and Restarted Convergence of PDUOT\\\large Working Theory Note}
\author{}
\date{October 12, 2026}

\begin{document}
\maketitle

\begin{abstract}
This note isolates a convergence result that is both valid for restricted KL-UOT and relevant to the implemented PDUOT solver.
For fixed step sizes, the PDUOT update is a reflected Halpern iteration applied to a firmly nonexpansive PDHG map.
This observation gives global convergence and an $O(1/k)$ fixed-point residual bound by existing operator theory.
The KL-specific contribution is a verifiable sufficient condition for local fixed-point sharpness: strict complementarity together with full column rank of the incidence matrix on the positive transport support.
For a bipartite transport graph, the rank condition holds when the positive support is a forest.
We prove strong metric regularity directly by identifying the inactive face and applying the implicit function theorem to the reduced KKT equations.
The resulting error bound yields tail linear convergence for fixed-frequency or residual-based restarts.
The result covers a fixed-metric core of PDUOT and deliberately excludes the PID primal-weight controller and the complete practical restart policy.
\end{abstract}

\section{Scope}

The theoretical target is one restricted KL-UOT problem on a fixed active support.
All convergence statements concern the exact-arithmetic proximal map and a certified step satisfying $\tau\sigma\lVert K\rVert_2^2<1$.
The proof has two layers.
The first layer identifies the fixed-metric PDUOT iteration with reflected Halpern PDHG and is largely inherited from existing theory \cite{lu2024restarted,liu2026restarted}.
The second layer establishes a local error bound from the geometry of KL-UOT.
Only the second layer is problem-specific.

The current implementation updates the primal weight at restarts with a PID controller and uses sufficient, necessary, and artificial restart triggers inherited from cuPDLPx \cite{lu2025cupdlpx}.
Those mechanisms are discussed separately in Section~\ref{sec:implementation-gap}.

\section{Restricted KL-UOT as a saddle problem}

Let $\mathcal A$ be the active support and let $x\in\R^{|\mathcal A|}$ contain its transport values.
Let $K\in\R^{(n_s+n_t)\times|\mathcal A|}$ be the unsigned bipartite incidence matrix, so that $Kx=(r,q)$ contains the source and target transported marginals.
For positive reference masses $a$ and $b$, restricted KL-UOT is
$$\min_{x\geq 0}\ c^\top x+\rho_a\sum_i\left(r_i\log\frac{r_i}{a_i}-r_i+a_i\right)+\rho_b\sum_j\left(q_j\log\frac{q_j}{b_j}-q_j+b_j\right),\qquad (r,q)=Kx.$$
Define $F(x)=c^\top x+\iota_{\R_+^{|\mathcal A|}}(x)$ and let $G(Kx)$ denote the two KL terms.
Its conjugate is
$$G^*(y)=\rho_a\sum_i a_i\left(e^{y_i/\rho_a}-1\right)+\rho_b\sum_j b_j\left(e^{y_{n_s+j}/\rho_b}-1\right).$$
The saddle formulation is $\min_x\max_y F(x)+\langle Kx,y\rangle-G^*(y)$.
The paper potentials satisfy $y=-(f,g)$.

The KKT mapping is
$$\mathcal M(x,y)=\begin{pmatrix}c+K^\top y+N_{\R_+^{|\mathcal A|}}(x)\\ \nabla G^*(y)-Kx\end{pmatrix}.$$
For steps $\tau,\sigma>0$, define the canonical PDHG metric
$$P=\begin{pmatrix}\tau^{-1}I&-K^\top\\-K&\sigma^{-1}I\end{pmatrix}.$$
The condition $\tau\sigma\lVert K\rVert_2^2<1$ makes $P$ positive definite.
A constant common step of $0.998/\lVert K\rVert_2$ satisfies this condition while preserving the product $\tau\sigma$.
The implementation uses a power-iteration estimate in place of the exact norm, so this statement becomes certified only when that estimate is replaced by a verified upper bound.

One PDHG image $T(x,y)=(\widehat x,\widehat y)$ is
$$\widehat x=\prox_{\tau F}(x-\tau K^\top y),\qquad \widehat y=\prox_{\sigma G^*}\left(y+\sigma K(2\widehat x-x)\right).$$
The exponential proximal map for $G^*$ is the Lambert--$W$ update implemented by PDUOT.

\begin{proposition}[Operator form]
If $P\succ0$, then $T=(P+\mathcal M)^{-1}P$ and $T$ is firmly nonexpansive in the $P$-norm.
Consequently, $2T-I$ is nonexpansive and $\Fix(T)=\mathcal M^{-1}(0)$.
\end{proposition}

\begin{proof}
The inclusion $P(z-Tz)\in\mathcal M(Tz)$ expands into the two displayed proximal updates.
The mapping $\mathcal M$ is maximally monotone because it is the KKT operator of a proper closed convex saddle problem.
The preconditioned resolvent is therefore firmly nonexpansive in the $P$-metric.
\end{proof}

Within one restart epoch, PDUOT applies the Halpern iteration to the reflected map:
$$z^{k+1}=\frac{1}{k+2}z^0+\frac{k+1}{k+2}(2T-I)z^k.$$

\begin{corollary}[Fixed-metric global convergence]
Suppose the restricted KL-UOT problem has a KKT point and $P\succ0$ is fixed.
Then the unrestarted reflected Halpern iterates converge to a fixed point of $T$ and satisfy
$$\lVert z^k-Tz^k\rVert_P\leq\frac{\dist_P(z^0,\Fix T)}{k+1}.$$
\end{corollary}

This corollary is a direct use of reflected Halpern fixed-point theory and should not be presented as a new KL-UOT theorem.

\section{A KL-specific local error bound}

Let $(x^*,y^*)$ be a KKT point.
Define the positive support $B=\{e:x_e^*>0\}$ and its complement $N=\mathcal A\setminus B$.
The dual slack is $s^*=c+K^\top y^*=c-f-g$.

\begin{lemma}[Positive optimal marginals]
If every source and target node is incident to at least one active-support edge, then every optimal transported marginal is strictly positive.
\end{lemma}

\begin{proof}
Let $x^*$ be optimal and suppose that $r_i^*=0$ for some source node $i$.
Choose an active-support edge $e=(i,j)$ and set $x(t)=x^*+t\mathbf e_e$ for $t>0$.
The source KL change divided by $t$ is $\rho_a(\log(t/a_i)-1)$, which tends to $-\infty$ as $t\downarrow0$.
The linear-cost quotient is the finite number $c_e$.
If $q_j^*>0$, the target KL quotient converges to the finite derivative $\rho_b\log(q_j^*/b_j)$; if $q_j^*=0$, it also tends to $-\infty$.
Hence the total directional quotient is negative for all sufficiently small $t$, contradicting optimality.
The target-side argument is identical.
\end{proof}

The lemma also gives existence of a finite KKT dual variable.
Indeed, the objective is differentiable at every optimal transported marginal, so the convex Fermat rule yields the stationarity inclusion with $y^*=\nabla G(Kx^*)$.
At a KKT point, $y^*=\nabla G(Kx^*)$ and $Kx^*=\nabla G^*(y^*)$.
The conjugate is locally strongly convex at $y^*$, although its smallest curvature can be arbitrarily small.
More precisely,
$$H:=\nabla^2G^*(y^*)=\operatorname{diag}\left(\frac{r^*}{\rho_a},\frac{q^*}{\rho_b}\right)\succ0.$$

\begin{assumption}[Nondegenerate positive support]\label{ass:nondegenerate}
The KKT point satisfies strict complementarity $s_e^*>0$ for every $e\in N$, and $K_B$ has full column rank.
\end{assumption}

\begin{lemma}[Graph interpretation]
The unsigned bipartite incidence matrix $K_B$ has full column rank if and only if the graph formed by the positive transport edges is a forest.
\end{lemma}

\begin{proof}
Multiplying all target rows by $-1$ turns $K_B$ into an oriented incidence matrix without changing its column rank.
The columns of an oriented incidence matrix are independent exactly when the corresponding edge set contains no cycle.
\end{proof}

\begin{lemma}[Reduced Jacobian]\label{lem:reduced-jacobian}
Under Assumption~\ref{ass:nondegenerate}, the matrix
$$J=\begin{pmatrix}0&K_B^\top\\-K_B&H\end{pmatrix}$$
is nonsingular.
\end{lemma}

\begin{proof}
If $J(\delta x,\delta y)=0$, then $K_B^\top\delta y=0$ and $-K_B\delta x+H\delta y=0$.
Taking the inner product of the second equality with $\delta y$ gives $\delta y^\top H\delta y=0$, hence $\delta y=0$.
Full column rank of $K_B$ then gives $\delta x=0$.
\end{proof}

\begin{theorem}[Local strong regularity]\label{thm:strong-regularity}
Under Assumption~\ref{ass:nondegenerate}, the inverse KKT mapping $\mathcal M^{-1}$ has a single-valued Lipschitz localization around $(0,z^*)$, where $z^*=(x^*,y^*)$.
Consequently, there are a neighborhood $U$ of $z^*$ and a constant $\kappa>0$ such that
$$\lVert z-z^*\rVert_2\leq\kappa\dist_2(0,\mathcal M(z))\qquad\text{for all }z\in U.$$
\end{theorem}

\begin{proof}
Write a perturbation of the KKT inclusion as $w=(u,v)\in\mathcal M(x,y)$.
For every positive edge $e\in B$, continuity preserves $x_e>0$ near $x^*$, so the normal cone vanishes and $u_B=c_B+K_B^\top y$.
When $N$ is nonempty, let $\delta=\min_{e\in N}s_e^*>0$ and shrink the neighborhoods so that $c_N+K_N^\top y\geq 2\delta/3$ and $\lVert u_N\rVert_\infty<\delta/3$.
If a nearby solution had $x_e>0$ for an edge $e\in N$, then the normal cone would vanish and $u_e=c_e+K_e^\top y\geq2\delta/3$, contradicting $|u_e|<\delta/3$.
Thus every nearby solution satisfies $x_N=0$; when $N$ is empty, this conclusion is vacuous.

The remaining equations are
$$c_B+K_B^\top y=u_B,\qquad \nabla G^*(y)-K_Bx_B=v.$$
Their Jacobian with respect to $(x_B,y)$ at $(x_B^*,y^*)$ is the nonsingular matrix in Lemma~\ref{lem:reduced-jacobian}.
The classical implicit function theorem therefore gives a unique continuously differentiable, hence locally Lipschitz, solution $(x_B(u_B,v),y(u_B,v))$ for all sufficiently small $(u_B,v)$.
For $e\in N$, the inclusion at $x_e=0$ is equivalent to $u_e\leq c_e+K_e^\top y$, which holds automatically in the chosen neighborhoods.
The localized solution is consequently independent of $u_N$ and is unique for every sufficiently small full perturbation $(u,v)$.
This proves strong metric regularity in the sense of \cite{robinson1980strongly}.
Strong metric regularity implies the displayed strong metric subregularity estimate.
\end{proof}

The proof also gives a local conditioning interpretation.
The derivative of the localized inverse at zero contains $J^{-1}$, and the Schur complement for the primal block is $K_B^\top H^{-1}K_B$.
Thus the error-bound constant is controlled by the KL marginal curvatures and the singular geometry of the positive-support incidence matrix.
A nearly vanishing entry of $H$ or a poorly conditioned $K_B$ can make the bound weak.

\begin{proposition}[Cycle obstruction]\label{prop:cycle-obstruction}
If the positive-support graph contains a cycle, then $x^*$ is not an isolated primal optimum.
Therefore the pointwise strong metric regularity conclusion of Theorem~\ref{thm:strong-regularity} cannot hold at $(x^*,y^*)$.
\end{proposition}

\begin{proof}
A cycle gives a nonzero vector $d_B\in\ker K_B$.
Because $x_B^*>0$, both $x_B^*+td_B$ and $x_B^*-td_B$ remain positive for sufficiently small $t$.
The transported marginals are unchanged.
Stationarity on $B$ gives $c_B=-K_B^\top y^*$ and hence $c_B^\top d_B=0$.
The two perturbed plans therefore have the same objective value as $x^*$.
\end{proof}

Proposition~\ref{prop:cycle-obstruction} does not rule out a weaker error bound to the full solution set.
It shows that a different argument would be needed for cyclic optimal supports.

\begin{corollary}[Fixed-point sharpness]\label{cor:sharpness}
Under the conditions of Theorem~\ref{thm:strong-regularity}, there are a neighborhood $V$ of $z^*$ and a constant $\alpha>0$ such that
$$\alpha\dist_P(z,\Fix T)\leq\lVert z-Tz\rVert_P\qquad\text{for all }z\in V.$$
One valid choice is
$$\alpha=\left(1+\kappa\lambda_{\max}(P)\sqrt{\frac{\lambda_{\max}(P)}{\lambda_{\min}(P)}}\right)^{-1}.$$
\end{corollary}

\begin{proof}
Let $\widehat z=Tz$.
The resolvent identity gives $P(z-\widehat z)\in\mathcal M(\widehat z)$.
Continuity of $T$ permits $V$ to be chosen so that $\widehat z\in U$ whenever $z\in V$.
Theorem~\ref{thm:strong-regularity} and the spectral bounds for $P$ give
$$\lVert\widehat z-z^*\rVert_2\leq\kappa\lVert P(z-\widehat z)\rVert_2\leq\frac{\kappa\lambda_{\max}(P)}{\sqrt{\lambda_{\min}(P)}}\lVert z-\widehat z\rVert_P.$$
The triangle inequality then yields
$$\dist_P(z,\Fix T)\leq\lVert z-z^*\rVert_P\leq\left(1+\kappa\lambda_{\max}(P)\sqrt{\frac{\lambda_{\max}(P)}{\lambda_{\min}(P)}}\right)\lVert z-Tz\rVert_P.$$
\end{proof}

\begin{theorem}[Tail linear convergence under certified restart]\label{thm:restart}
Suppose Assumption~\ref{ass:nondegenerate} holds and $P$ is fixed.
There is a neighborhood $W$ of $z^*$ with the following properties.
For any restart anchor $a^0\in W$, fixed-frequency reflected Halpern PDHG with epoch length $L$ satisfying $L+1\geq e/\alpha$ obeys
$$\dist_P(a^{n+1},\Fix T)\leq e^{-1}\dist_P(a^n,\Fix T).$$
Alternatively, terminate each epoch at the first $k$ satisfying
$$\lVert z^k-Tz^k\rVert_P\leq\theta\lVert a^n-Ta^n\rVert_P$$
for a fixed $\theta\in(0,1)$ and restart from $a^{n+1}=Tz^k$.
Then every epoch is finite and
$$\lVert a^n-Ta^n\rVert_P\leq\theta^n\lVert a^0-Ta^0\rVert_P,\qquad \dist_P(a^n,\Fix T)\leq\alpha^{-1}\theta^n\lVert a^0-Ta^0\rVert_P.$$
\end{theorem}

\begin{proof}
Choose a closed $P$-ball around $z^*$ contained in the sharpness neighborhood $V$ and let $W$ be its interior.
Both $T$ and $2T-I$ are nonexpansive and fix $z^*$.
The Halpern convex combination therefore keeps every inner iterate and every PDHG image in the ball whenever its anchor lies in $W$.

For an epoch anchored at $a$, the reflected Halpern residual estimate is
$$\lVert z^k-Tz^k\rVert_P\leq\frac{\dist_P(a,\Fix T)}{k+1}.$$
At a restart $a^+=Tz^k$, nonexpansiveness of $T$ gives
$$\lVert a^+-Ta^+\rVert_P=\lVert Tz^k-T(Tz^k)\rVert_P\leq\lVert z^k-Tz^k\rVert_P.$$
Corollary~\ref{cor:sharpness} at $a^+$ consequently gives
$$\dist_P(a^+,\Fix T)\leq\frac{1}{\alpha(k+1)}\dist_P(a,\Fix T).$$
The fixed-frequency conclusion follows when $k=L$ and $L+1\geq e/\alpha$.

For residual-based restart, the second displayed inequality and the restart condition imply $\lVert a^{n+1}-Ta^{n+1}\rVert_P\leq\theta\lVert a^n-Ta^n\rVert_P$.
Induction gives the residual bound, and Corollary~\ref{cor:sharpness} gives the distance bound.
Finally, sharpness at the anchor and the reflected Halpern residual estimate imply
$$\lVert z^k-Tz^k\rVert_P\leq\frac{1}{\alpha(k+1)}\lVert a-Ta\rVert_P,$$
so the restart condition must hold once $k+1\geq1/(\alpha\theta)$.
\end{proof}

\section{Relation to existing theory}

The operator form, global convergence, and reflected Halpern residual rate are standard consequences of fixed-point theory.
Recent work also gives restarted reflected Halpern guarantees for broad composite primal--dual methods under an assumed local error bound \cite{liu2026restarted}.
Repeating those abstract results for PDUOT has little independent value.

The useful result is Theorem~\ref{thm:strong-regularity} and its graph interpretation.
It replaces an abstract sharpness assumption by conditions expressed in transport quantities: inactive-edge slack, positive marginals, and the cycle structure of the positive support.
This is suitable as a supporting theorem in a HierUOT paper.
It is unlikely to support a separate theory paper without a sharper constant, a weaker condition that allows cyclic optimal supports, or a global result.

\section{Implementation gap}\label{sec:implementation-gap}

\begin{table}[ht]
\centering
\caption{Coverage of the current PDUOT implementation by the proposed proof.}
\begin{tabular}{lll}
\toprule
Component & Current implementation & Covered here \\
\midrule
PDHG map and KL proximal step & Yes & Yes \\
Reflection and Halpern anchoring & Yes & Yes \\
Constant common step & Default & Yes \\
Fixed primal weight & Optional & Yes \\
Residual reduction by $e^{-1}$ & Available as a variant & Yes \\
$0.2/0.5/0.36$ practical restart rules & Default & No \\
PID primal-weight update & Default & No \\
Adaptive common step & Optional & No \\
Six-step Lambert--$W$ evaluation & Default & No \\
Power-iteration norm estimate & Default & Needs a certified upper bound \\
\bottomrule
\end{tabular}
\end{table}

The sufficient restart trigger with factor $0.2$ is compatible with a residual contraction argument when it fires.
The necessary trigger allows factor $0.5$, and the artificial trigger does not require residual reduction.
They therefore do not coincide with the certified $e^{-1}$ rule.

The PID controller changes $\tau$ and $\sigma$ across epochs while preserving their product.
Each epoch still has a valid PDHG metric, but the Fej\'er and sharpness arguments use different norms after successive restarts.
A variable-metric extension would require explicit control of metric changes, such as uniform equivalence plus a summable relative-change condition.
The present PID safeguards bound the weight but do not impose summable changes.
The full-controller proof is therefore a separate problem rather than a direct corollary.
The CUDA proximal kernel also uses a fixed number of Newton steps to evaluate the Lambert--$W$ expression, and the spectral norm is estimated by power iteration.
These are floating-point realizations of the exact map analyzed above.
A code-level convergence theorem would additionally require a controlled inexact-prox sequence and a certified upper bound on $\lVert K\rVert_2$.

\section{Diagnostic evidence}

The following experiments test whether the sufficient conditions occur in representative solutions.
They are diagnostics, not a convergence study.

% Data source: papers/hier_uot/pduot_convergence_diagnostics.json, geometry_experiment.
\begin{table}[ht]
\centering
\caption{Local geometry of high-accuracy HierUOT solutions on TOME with $n_s=n_t=2048$, FP64 costs, and tolerance $10^{-8}$.}
\begin{tabular}{rrrrr}
\toprule
$\rho_a=\rho_b$ & Positive edges & Forest & Covered nodes & Minimum inactive slack \\
\midrule
$1$ & 3982 & yes & $2048/2048$ & $4.76\times10^{-6}$ \\
$0.1$ & 3674 & yes & $2048/2048$ & $6.22\times10^{-6}$ \\
$0.01$ & 3214 & yes & $2041/2024$ & $1.04\times10^{-6}$ \\
\bottomrule
\end{tabular}
\end{table}

The first two cases satisfy the numerical counterparts of forest support, strict complementarity, and positive node coverage.
At $\rho=0.01$, several computed marginals reach numerical zero even though exact optimal marginals are positive under node coverage.
The theorem remains qualitatively applicable to the exact problem, but its curvature-dependent constant is unlikely to predict the observed finite-precision behavior in this regime.

% Data source: papers/hier_uot/pduot_convergence_diagnostics.json, restart_experiment.
\begin{table}[ht]
\centering
\caption{Cold restricted solves on the 1958 positive edges of one TOME solution with $n_s=n_t=1024$, $\rho_a=\rho_b=1$, and tolerance $10^{-8}$.}
\begin{tabular}{lrr}
\toprule
Variant & Inner iterations & Restarts \\
\midrule
Fixed metric, certified restart & 27,000 & 11 \\
Fixed metric, practical restart & 28,000 & 11 \\
PID metric, practical restart & 4,600 & 8 \\
\bottomrule
\end{tabular}
\end{table}

The fixed-metric certified variant converges and displays geometric residual reduction across restart anchors, as predicted.
The PID controller is substantially faster on this instance.
This gap is a reason to keep the controller in the implementation while stating the theorem for a controlled core variant.

\section{Recommended paper claim}

The main paper can state three levels of results.
First, fixed-metric PDUOT is reflected Halpern PDHG and inherits global convergence with an $O(1/k)$ fixed-point residual.
Second, strict complementarity and a forest positive support imply a KL-specific local error bound.
Third, a certified restart variant then has tail linear convergence.

The practical PID and restart rules should be described as computational enhancements shared with cuPDLPx.
Their empirical value can be shown by ablation, while the theorem explicitly applies to the fixed-metric certified variant.

\begin{thebibliography}{9}

\bibitem{lu2024restarted}
H.~Lu and J.~Yang.
\newblock Restarted Halpern PDHG for Linear Programming.
\newblock \emph{arXiv preprint arXiv:2407.16144}, 2024.

\bibitem{lu2025cupdlpx}
H.~Lu, Z.~Peng, and J.~Yang.
\newblock cuPDLPx: A Further Enhanced GPU-Based First-Order Solver for Linear Programming.
\newblock \emph{arXiv preprint arXiv:2507.14051}, 2025.

\bibitem{liu2026restarted}
B.~Liu, J.~Cao, W.~Yin, and Z.~Wen.
\newblock Restarted Reflected Halpern Acceleration for Augmented Primal--Dual Methods.
\newblock \emph{arXiv preprint arXiv:2606.16552}, 2026.

\bibitem{robinson1980strongly}
S.~M. Robinson.
\newblock Strongly regular generalized equations.
\newblock \emph{Mathematics of Operations Research}, 5(1):43--62, 1980.

\end{thebibliography}

\end{document}
